\documentclass[11pt]{article}
\usepackage{amsmath,amssymb,amsthm}
\usepackage[margin=1.2in]{geometry}
\newtheorem{theorem}{Theorem}
\newtheorem{lemma}[theorem]{Lemma}
\newtheorem{remark}[theorem]{Remark}
\newcommand{\fpart}[1]{\{#1\}}
\DeclareMathOperator{\dimH}{dim_H}

\title{Bounded reciprocal mass of small non-divisor primes of $\binom{2n}{n}$,\\
outside a power-saving exceptional set}
\author{Samuel Lavery}
\date{Draft --- \today}

\begin{document}
\maketitle

\begin{abstract}
Call a prime $p$ \emph{balanced at $n$} if $p\nmid\binom{2n}{n}$; by Kummer's theorem
this holds iff every base-$p$ digit of $n$ is at most $(p-1)/2$.  For a constant
cutoff $B_1$ we prove: for every $\varepsilon>0$ and all large $N$, all but at most
$N^{\varepsilon}$ integers $n\in[N,2N]$ satisfy
\[
\sum_{\substack{3\le p\le B_1\\ p\nmid\binom{2n}{n}}}\frac1p\;\le\;
\mathrm{LP}(B_1)+\varepsilon,
\qquad
\mathrm{LP}(B_1):=\max\Big\{\sum_{p\in S}\tfrac1p:\ S\subseteq\{3,\dots,B_1\}
\text{ prime},\ \sum_{p\in S}\beta_p\le1\Big\},
\]
where $\beta_p=1-\log\frac{p+1}{2}/\log p$.  One has $\mathrm{LP}(B_1)\le
\mathrm{LP}(\infty)=0.6941\ldots$, against the trivial bound $\sum_{p\le B_1}1/p\to\infty$.
The proof identifies joint balance with membership of the mantissa in an intersection of
self-similar Cantor sets and applies the higher-rank Furstenberg slicing theorem of
Corso--Shmerkin (arXiv:2409.04608, Corollary 1.17), a preprint whose relevant proof chain
we have verified independently; the result is conditional on it and on nothing else
beyond elementary arithmetic.  Empirically the exceptional set appears to be not merely
power-saving but finite: the integers balanced at $\{3,5,7,11\}$ simultaneously are
$\{1,3160\}$ up to $10^{9}$.
\end{abstract}

\section{Statement}

Fix once and for all an odd prime bound $B_1$ and write $\mathcal P=\{3\le p\le B_1:
p\text{ prime}\}$.  For an odd prime $p$ let
\[
D_p=\{0,1,\dots,\tfrac{p-1}2\},\qquad
C_p=\Big\{\sum_{j\ge1}d_jp^{-j}:d_j\in D_p\Big\}\subset[0,1],
\]
the \emph{lower-half-digit Cantor set}, a homogeneous self-similar set with contraction
ratio $1/p$, digit set $D_p$, satisfying the open set condition, with
\[
\dimH C_p=\frac{\log\frac{p+1}2}{\log p}=1-\beta_p,\qquad
\beta_p:=\frac{\log\frac{2p}{p+1}}{\log p}.
\]
Say $n$ is \emph{$p$-balanced} if every base-$p$ digit of $n$ lies in $D_p$; by Kummer's
theorem this is equivalent to $p\nmid\binom{2n}{n}$.  Let $S(n)=\{p\in\mathcal P:\ n\ \text{is}\
p\text{-balanced}\}$.

\begin{theorem}\label{thm:core}
Assume Corollary 1.17 of arXiv:2409.04608 \emph{(see Remark~\ref{rem:dependency})}.
For every $\varepsilon>0$ there are $C_0=C_0(B_1,\varepsilon)$ and $N_0$ such that for all
$N\ge N_0$,
\[
\#\Big\{n\in[N,2N]:\ \sum_{p\in S(n)}\frac1p>\mathrm{LP}(B_1)\Big\}\;\le\;C_0\,N^{\varepsilon}.
\]
\end{theorem}

Since $\mathrm{LP}(B_1)\le\mathrm{LP}(\infty)<0.695$ while $\sum_{p\in\mathcal P}1/p$ is
unbounded in $B_1$, the theorem separates the balanced mass from its trivial bound by an
unbounded margin, for all but a power-saving number of integers.

\section{Proof}

\subsection{Over-budget sets are Cantor intersections}
If $\sum_{p\in S(n)}1/p>\mathrm{LP}(B_1)$ then, by definition of the linear program,
$\sum_{p\in S(n)}\beta_p>1$; discarding elements, $S(n)$ contains a subset $S$ with
\begin{equation}\label{eq:overbudget}
\sum_{p\in S}\beta_p>1,\qquad\text{i.e.}\qquad
s:=\sum_{p\in S}\dimH C_p<|S|-1=:k-1 .
\end{equation}
It therefore suffices to bound, for each of the at most $2^{|\mathcal P|}$ (a constant)
subsets $S$ satisfying \eqref{eq:overbudget}, the number of $n\in[N,2N]$ balanced at every
$p\in S$.

Fix such an $S=\{p_1,\dots,p_k\}$ and let $n\in[N,2N]$ be balanced at each $p_j$.  Write
$d_j=\lfloor\log 2N/\log p_j\rfloor$, so that $p_j^{d_j}\le 2N<p_j^{d_j+1}$ and $n$ has at
most $d_j$ base-$p_j$ digits.  Balance says exactly that
\begin{equation}\label{eq:membership}
\frac{n}{p_j^{d_j}}\;\in\;C_p^{(d_j)}\qquad(1\le j\le k),
\end{equation}
where $C_p^{(J)}$ denotes the depth-$J$ truncation of $C_p$ (the union of the $|D_p|^J$
level-$J$ construction intervals).

\subsection{The warp: mantissa coordinates}\label{sec:warp}
Rescale by the window: put $x=n/(2N)\in[\tfrac12,1]$.  Then \eqref{eq:membership} reads
\[
x\in g_j\big(C_{p_j}^{(d_j)}\big),\qquad
g_j(t)=\frac{p_j^{d_j}}{2N}\,t,
\qquad\text{with slopes}\qquad
\frac{1}{p_j}<\frac{p_j^{d_j}}{2N}\le 1 .
\]
This is the decisive normalization: \emph{the affine maps have slopes bounded between
$1/B_1$ and $1$, independently of $N$.}  (Applied in value space the slopes would be of
size $N$ and the uniform constant below would be vacuous.)

\subsection{Truncation}
The sets in \S\ref{sec:warp} are truncations, not the invariant sets themselves.  The
bridge is:

\begin{lemma}[truncation bridge]\label{lem:T}
Let $C_1,\dots,C_k$ be the lower-half Cantor sets of distinct odd primes, $s$ as above,
$\varepsilon>0$.  There is $C^\star(k,\varepsilon,p_1,\dots,p_k)$ such that for all depths
$J_j\ge1$ and all $\delta$ with $\max_jp_j^{-J_j}\le\delta\le\delta_0$,
\[
N_\delta\Big(\bigcap_{j=1}^k g_j\big(C_j^{(J_j)}\big)\Big)\;\le\;
7\,(3\sqrt k)^{k-1}\,C^\star\,\delta^{-\max\{s-(k-1),0\}-\varepsilon},
\]
uniformly over affine $g_j$ with slopes in $[1/B_1,1]$.
\end{lemma}

\begin{proof}
Each point of a depth-$J$ truncation lies within $p_j^{-J_j}\le\delta$ of $C_j$ (append
zero digits; $0\in D_p$).  A $3\delta$-separated subset of the intersection therefore
lifts, coordinatewise, to a $\delta$-separated subset of
$(g_1(C_1)\times\cdots\times g_k(C_k))\cap\pi^{-1}(B(0,\delta\sqrt k))$, where $\pi$ is the
orthogonal projection killing the diagonal.  Covering $B(0,\delta\sqrt k)$ by
$(3\sqrt k)^{k-1}$ $\delta$-balls and applying the ball-fibre bound underlying Corollary
1.17 of arXiv:2409.04608 (their Lemma 6.1, with the slopes absorbed into the constant via
$K=B_1$) bounds each fibre by $C^\star\delta^{-\max\{s-(k-1),0\}-\varepsilon}$; a
$3\delta$-packing/$\delta$-covering comparison on $\mathbb R$ costs the factor $7$.
\end{proof}

\subsection{Counting and union}
Apply Lemma~\ref{lem:T} with $\delta=1/(2N)$ and $J_j=d_j$ (so
$p_j^{-d_j}<p_j/(2N)\le B_1\delta$; replace $\delta$ by $B_1\delta$, absorbed in
constants).  By \eqref{eq:overbudget}, $\max\{s-(k-1),0\}=0$, so the number of
$\delta$-cells of $[\tfrac12,1]$ meeting the intersection is at most
$C_1(B_1,\varepsilon)\,N^{\varepsilon}$.  Each $\delta$-cell has length $1/(2N)$ and thus
contains at most one value of $x=n/(2N)$ with $n$ integral.  Hence at most
$C_1N^{\varepsilon}$ integers $n\in[N,2N]$ are balanced at all of $S$.  Summing over the
constantly many over-budget $S$ gives the theorem with
$C_0=2^{\pi(B_1)}\max_SC_1$. \qed

\section{Remarks}

\begin{remark}[dependency]\label{rem:dependency}
The input is Corollary 1.17 of E.~Corso and P.~Shmerkin, \emph{Dynamical self-similarity,
$L^q$-dimensions and Furstenberg slicing in $\mathbb R^d$}, arXiv:2409.04608 (2024), a
preprint at the time of writing.  We have verified the proof chain we consume
(Cor.~1.17 $\Leftarrow$ Thm.~1.15 $\Leftarrow$ Thm.~5.4 $\Leftarrow$ their inverse-theorem
engine) for the special case used here, in the course of which two repairable defects were
identified (a misprinted hypothesis (P1) of their \S5, whose intended form is confirmed by
their Remark 1.16, and an omitted almost-everywhere injectivity verification in the last
step of Theorem 5.4, which we supply for our projection); details in the project's
verification ledger.  No other non-elementary input occurs.
\end{remark}

\begin{remark}[what is and is not proved]
The theorem bounds the balanced mass of primes up to a \emph{constant} $B_1$, for all $n$
outside a power-saving exceptional set.  It does not bound the mass over the growing range
$p\le(\log n)^{O(1)}$ (which requires a base-uniform, effective form of the slicing
theorem), and it does not empty the exceptional set (a rigidity question).  Empirically
the exceptional sets appear finite and tiny: the integers balanced simultaneously at
$\{3,5,7,11\}$ --- the smallest over-budget set, $\sum\beta=1.227$ --- are $\{1,3160\}$
up to $10^9$, and at $\{3,5,7,13\}$ are $\{1,756,757,3250\}$ up to $10^8$.
\end{remark}

\begin{remark}[context]
$\mathrm{LP}(\infty)=0.6941\ldots$ matches, to the accuracy of measurement, the empirical
ceiling of the balanced mass located by exhaustive search in earlier stages of this
programme ($\approx0.69$--$0.83$ across normalizations), suggesting the linear program is
the true extremal mechanism.  The theorem is a step in a programme to prove
$\sum_{p\nmid\binom{2n}{n}}1/p=O(1)$ for all $n$ (the bounded form of Erd\H os problem
377); the asymptotic form is treated in a companion paper.
\end{remark}

\end{document}
